\scanpage{43}
% archive sheet number: 42
essentiels sont

a) Les $\unsure{\nu|v}(\Pi^-_{f_2})$ --- \unsure{dénoté} $\unsure{p|\rho}_0(\Pi)$ si $\Pi$ \uline{régulier}

b) les $\nu(\Pi(d_{0,3})),\nu(\Pi(d_{1,4}))\cdots(\nu(\Pi(d_{n-3,n})))$

notés $p_1(\Pi),p_2(\Pi),\ldots,p_{n-2}(\Pi)$ si $\Pi$ \uline{régulier}

c) les $\nu(\Pi^+_{f_{n-2}})$, \unsure{dénoté|noté} $p_{n-1}(\Pi)$ si $\Pi$ \uline{régulier}

\par\noindent\hspace*{.25\linewidth}\rule{.25\linewidth}{.4pt}\par

Supposons $\Pi$ \uline{régulier}, je dis que
\[
\left\{\begin{aligned}
p_0&=\text{ordre de }\sigma_0\sigma_1\\
p_1&:\text{ordre de }\sigma_1\sigma_2\\
&\cdots\\
p_{n-1}&:\text{Ordre de }\sigma_{n-1}\sigma_n
\end{aligned}\right.
\]
\margin{vrai tout au moins dans le cas des géométries d'incidences (``polyèdres''\ldots)}

En effet, \deleted{\unsure{$\Pi$}} sous hyp. de \uline{régularité}, $\nu(\Pi^-_{f_2})$ est l'ordre de $\sigma_0\sigma_1$ opérant sur $E^-(2)_{f_2}\inserted{=(E/G_n^+(2))_{f_2}}$, \deleted{\unsure{$\nu(\Pi(d^{\unsure{a|o}}_{i,i+3}))$ est l'ordre}}
\[
(\sigma_0\sigma_1)^p x=\deleted{\unsure{x}}g_x\deleted{\unsure{\cdot}}x\qquad g_x\in G_n^+(2)
\]
\diagram{\(\begin{array}{ccccc}
i-1&&\unsure{i+1\leq j-1}&j+1&\\
&i&j&&i\\
&&\unsure{\geq|=}&&\\
&&i+2&&
\end{array}\)}
